\documentclass[11pt]{article}
\usepackage[a4paper,margin=26mm]{geometry}
\usepackage{amsmath,amssymb}
\usepackage[hidelinks]{hyperref}
\setlength{\emergencystretch}{3em}
\title{Conjecture 00000002894: a vertexless triangle in $\mathrm{PG}(2,9)$}
\author{lizaixi01, with OpenAI Codex assistance}
\date{23 October 2026}
\begin{document}
\maketitle
\section{The classification fails without a smallness restriction}
The English statement claims that nontrivial minimal blocking sets in
planes of order $q^2$ are of R\'edei type, and that non-R\'edei examples
exist only when $q$ is square. Neither clause imposes a smallness
restriction. We give an inclusion-minimal counterexample in the
Desarguesian plane of order $Q=9=3^2$. Its parameter $q=3$ is not square.
Here ``minimal'' means inclusion-minimal, as in the English statement
and standard terminology; it does not mean minimum cardinality.

A blocking set meets every line. It is nontrivial if it contains no
whole line, and minimal if no proper subset is blocking. A R\'edei
blocking set has a line $\ell$ with $|B\setminus\ell|=Q$.
These definitions appear in G\'acs--Lov\'asz--Sz\H{o}nyi,
\href{https://msp.org/iig/2008/6-1/iig-v6-n1-p12-p.pdf}
{\emph{Directions in $\mathrm{AG}(2,p^2)$}},
Innovations in Incidence Geometry 6--7 (2007--08), pp.189--201,
printed p.191. In particular, the necessary property of R\'edei type
used here does not require an additional size bound.

\section{The actual point set and its tangents}
Work over $K=\mathbf F_9$, in homogeneous coordinates. Let
\[
 B=\{[0:1:t],[t:0:1],[1:t:0]:t\in K^\times\}.
\]
These are precisely the three coordinate sides with their three
vertices removed. The three families are disjoint, and each has
$8$ distinct points, so $|B|=24$.

Every nonzero linear form $a_0x+a_1y+a_2z$ vanishes at some point of
$B$. If two coefficients are nonzero, choose the corresponding
coordinates to cancel and set the third coordinate to zero. If just
one coefficient is nonzero, choose its coordinate to be zero and
the other two to be one. Thus $B$ is blocking.

The tangent to $[0:1:t]$ has equation $-ty+z=0$; cyclically the other
tangents are $x-tz=0$ and $-tx+y=0$. Each meets $B$ in exactly the
indicated point. Removing any point therefore leaves a line unblocked.
More generally, any blocking subset of $B$ must contain every point
of $B$, which proves inclusion-minimality.

\newpage
\section{Nontriviality and failure of R\'edei type}
Every line has a point outside $B$. If a coefficient $a_i$ of its
nonzero linear form is zero, the corresponding coordinate vertex
lies on the line and is outside $B$. Otherwise the point
\[
 [a_0^{-1}:a_1^{-1}:a_2^{-1}]
\]
lies on the line because $1+1+1=0$ in characteristic three. All its
coordinates are nonzero, so it is outside $B$. Hence $B$ contains no
line and is nontrivial.

A projective line over $K$ has $Q+1=10$ points. Consequently every
line satisfies
\[
 |B\setminus\ell|=24-|B\cap\ell|\ \ge 24-10=14>9.
\]
There is no R\'edei line. Thus this same actual set refutes both the
universal R\'edei assertion and the claimed square-$q$ restriction.
The further characterization of R\'edei sets cannot repair either
false necessary clause and is not needed for the disproof.

\section{Formal correspondence}
The Lean point type is Mathlib's genuine
\texttt{Projectivization K (Fin 3 $\to$ K)}, not a list of assumed
incidence objects. Lines are the projectivizations of kernels of
all nonzero homogeneous linear forms. The kernel dimension is
proved to be two by surjectivity of the form and rank--nullity.
A restriction map into the projectivization of that kernel is
proved injective, giving the line cardinality upper bound from
Mathlib's finite projective-space cardinality theorem. Only this
upper bound is needed in the formal contradiction.

The point set is the image of the explicit three families, and the
family map is proved injective using projective scalar equivalence.
Blocking, the exact tangent intersections, minimality and
nontriviality are proved from coordinates. The formal predicate
\texttt{IsRedei} is exactly the displayed necessary line-complement
criterion. The final \texttt{counterexample} specializes to
\texttt{GaloisField 3 2}, proves its cardinality is $3^2=9$, proves
$3$ is not square, and bundles the $24$-point count and all four
blocking/minimal/nontrivial/non-R\'edei properties. No stochastic
or numerical incidence test replaces the mathematical objects.

\section{Verification}
The project pins Lean 4.33.0 and Mathlib commit
\[
\texttt{db584cd6d46c92f209a44c0f1c829460d327499d}.
\]
Source-bound receipts accompany a new Main build, transitive axiom
audits and official \texttt{leanchecker --fresh Main} replay.
No \texttt{sorry}, \texttt{native\_decide} or additional axiom is used.
OpenAI Codex assisted; a separate supervising Codex agent reviewed
the formal chain. This is AI review, not human certification.
Organizer acceptance is pending.
\end{document}
